Three features of this match bear on its credibility. It is a match to an administrative registry
rather than to a survey or a self-report, so the information-status indicator does not depend on
seller cooperation and cannot be strategically misreported. It requires a geocoded address, so the
matched sample is the subset of listings that can be located at the address level rather than the
full listing universe; sample scopes are stated alongside the statistics reported below. And, most
consequentially, $K=0$ must be read as \emph{no post-July-2021 certificate in the public registry}
rather than as \emph{no certificate exists}. The registry was populated from the reform date
onward, while a certificate issued earlier could remain legally valid at the time of listing, so a
seller could hold and display a valid certificate that the registry cannot corroborate. The data
bear this out. The share of listings displaying a score with no registry record falls monotonically
as the registry accumulates, from $62.5\%$ of apartment listings in 2021Q3 to $18.4\%$ in 2024Q4,
mirrored by a rise in verified disclosers from $7.4\%$ to roughly one half.
